\begin{beginnerstatement}[Kurz erklärt: vom Versionsstand zum Paket]

Eine Version ist eine eindeutige Bezeichnung für einen bestimmten Stand; im
Template ist \texttt{VERSION} die zentrale Quelle dafür. Ein Release ist eine
zur Veröffentlichung freigegebene Version, doch der untersuchte Workflow
veröffentlicht selbst nichts. \texttt{release check} kontrolliert lediglich
die Voraussetzungen und erzeugt keine Artefakte. Beim \emph{Packaging} werden
Programmdateien zu einem auslieferbaren Paket oder \emph{Bundle}
zusammengefasst; dieses Ergebnis ist ein Artefakt. Ein Git-Tag markiert einen
bestimmten Commit mit einem Namen, häufig der Versionsnummer. Beim
\emph{Signing} bestätigt eine digitale Signatur mit einem Zertifikat die
Herkunft eines Pakets; Notarisierung ist eine zusätzliche externe Prüfung, und
ein Updater verteilt spätere Versionen an installierte Anwendungen. Deployment
bedeutet das Bereitstellen einer Anwendung; Docker baut dafür Images,
Container führen sie aus und Docker Compose beschreibt mehrere
zusammengehörige Container. AWS, Azure und GCP sind mögliche Cloud-Anbieter,
auf die sich diese anbieterneutrale Basis jedoch nicht automatisch festlegt.

\end{beginnerstatement}
